\section{Related Work}
\label{sec:related}

\paragraph{AI-generated content on the Web.}
The scale of machine-written text on the live Web has moved from speculation to measurement: as of mid-2026, an estimated 27.5\% of tokens passing common quality filters on the open Web are AI-generated, rising to 31.1\% within two months, and high-quality human text is becoming scarce enough that its token value can be modeled with scaling laws \cite{russell2026much}. Domain studies document the same shift operationally, from music streaming catalogs flooded with algorithmically produced tracks \cite{wu2026empirical} to scientific publishing \cite{oh2026science}, and a growing conceptual literature debates what ``slop'' is and why it matters \cite{kommers2025slop,nishal2026slop}. Human-factors work shows that both lay readers \cite{zhu2026humans} and automated pipelines struggle to distinguish such content. This literature treats synthetic content primarily as a \emph{data-quality} problem for training corpora and a \emph{perception} problem for readers. Our work complements it by asking a downstream question those studies do not touch: what happens when the consumers of this content are \emph{autonomous agents} that act on it.

\paragraph{Generative engine optimization and content manipulation.}
A second line of work makes the production side explicit. GEO methods rewrite pages to maximize citation by generative engines \cite{wu2025what}, and a rapidly growing benchmark suite measures the resulting arms race: E-GEO formalizes seller-side optimization in e-commerce \cite{bagga2025egeo}, GEO-Bench and its multimodal extension quantify ranking manipulation \cite{nimase2026geobench,du2026multimodal}, and IF-GEO studies multi-query optimization \cite{zhou2026ifgeo}. Detection responses are appearing at matching speed---GEO-Flag \cite{chu2026geoflag}, structural features that shape citation behavior \cite{yu2026structural}, and Counter-GEO-Bench \cite{zheng2026countergeobench} together with defense-oriented work \cite{li2026when}. Two properties of this literature shape our contribution. First, manipulation is almost always \emph{targeted}: an adversary crafts pages for specific queries or documents. Second, the affected system is the \emph{engine}: what gets distorted is retrieval ranking or generated answers. We study the complementary regime---ambient, non-targeted contamination generated for commercial motives at scale---and the affected system is a \emph{decision-making agent} that must commit to a purchase.

\paragraph{Agents under a polluted Web.}
Closest to us are systems that evaluate agents against corrupted information sources. The Synthetic Web builds programmatic mini-internets with adversarially curated pages and finds that language agents inherit epistemic weaknesses under single-injection falsehoods \cite{shah2026synthetic}. Endorsement-vulnerability measurements show that LLM search agents can be steered by attacker-published pages \cite{chen2026much}. A separate security thread demonstrates prompt-injection and memory-poisoning attacks on web agents \cite{cao2026ward,yang2026zombie,syros2026muzzle,guo2026logintrap}, and persuasion studies show agent behavior can be shifted by conversational or inter-agent influence \cite{jeong2026persuasion,qiu2026belief,zhou2026epistemic}, with Sybil-style analysis questioning what counts as independent evidence \cite{bara2026epistemic}. At the retrieval layer, Yu et al.\ show that cumulative AI pollution eventually collapses search quality \cite{yu2026retrieval}. All of these works attack the agent with \emph{deliberate, targeted} content, or measure degradation \emph{upstream} of the decision. None asks the question we pose: whether \emph{ordinary} optimization pressure---the same content farms that already dominate real e-commerce---degrades \emph{decision quality} in a dose-dependent way, and what class of defenses restores it. To our knowledge, no prior work measures agent choice-level susceptibility with a forced-choice probe, a decision-level dose--response function, and regret-based economics on a grounded commerce benchmark.

\paragraph{Shopping-agent benchmarks.}
WebShop \cite{yao2022webshop} established grounded language agents over real e-commerce data with attribute-level rewards; successors extend it to personalized curation \cite{kim2026agenticshop}, multi-shop comparison \cite{peeters2025webmall}, and deep research shopping \cite{lyu2025deepshop}. ReAct-style prompting \cite{yao2022react} remains the standard operating mode for these agents. We reuse WebShop's official data and reward function---unmodified---as the substrate for a controlled candidate-selection environment, which lets us attribute every reward difference to the manipulated variable (which listing the agent chose) rather than to benchmark artifacts.
